\documentclass[aspectratio=169]{beamer}
\usetheme{Madrid}
\usecolortheme{beaver}

\usepackage[UTF8,fontset=fandol]{ctex}
\usepackage{amsmath, amssymb}
\usepackage{booktabs}
\usepackage{ragged2e}
\usepackage{array}
\usepackage{tikz}
\usetikzlibrary{positioning, arrows.meta, shapes.geometric, calc, fit}

\newcommand{\framebottomnote}[1]{%
    \vfill
    \par\noindent
    {\tiny\color{black!65}\parbox{0.96\linewidth}{#1}\par}%
    \vspace*{-0.75cm}%
}

\newcolumntype{L}[1]{>{\RaggedRight\arraybackslash}p{#1}}

\title[Helplessness Update]{Helplessness Update 机制设计与文献对应}
\subtitle{Digital Friction 实验 --- 进展汇报}
\date{2026-04-08}

\begin{document}

\begin{frame}
    \titlepage
\end{frame}

% ============================================================
\section{背景与目标}
% ============================================================

\begin{frame}{旧版问题与改进目标}

\begin{columns}[T]
\begin{column}{0.48\textwidth}
\textbf{旧版（v0 baseline）存在的问题}
\begin{enumerate}
    \item \texttt{BASE\_DELTAS} 过大，outcome 直接主导 helplessness 变化
    \item 没有心理中介层（self-efficacy / felt control）
    \item 所有 avoid 一律计入 helplessness
    \item support 是万能直接减分器
    \item 无长期保护机制（controllable success memory）
    \item 无阻尼，反馈环容易数值爆炸
\end{enumerate}
\end{column}

\begin{column}{0.48\textwidth}
\textbf{改进后的设计原则}
\begin{enumerate}
    \item 主驱动改为 perceived uncontrollability
    \item 加入 self-efficacy / felt control 中介层
    \item avoid 拆分为 helpless / risk / low-value
    \item support 以间接修复（SE / CSM）为主
    \item 引入 controllable success memory 做乘性保护
    \item 加入 bounded-state damping
\end{enumerate}
\end{column}
\end{columns}

\framebottomnote{改进方向来自 learned helplessness / controllability 理论主线（Seligman 1967; Abramson 1978; Maier \& Seligman 2016）和老年数字服务场景文献。}
\end{frame}

% ============================================================
\section{机制总览}
% ============================================================

\begin{frame}{Helplessness Update --- 机制流程图}
\centering
\begin{tikzpicture}[
    node distance=0.55cm and 0.6cm,
    every node/.style={font=\scriptsize},
    box/.style={draw, rounded corners=3pt, minimum height=0.7cm, text width=2.1cm, align=center, fill=white},
    mainbox/.style={box, fill=blue!8, draw=blue!50},
    classbox/.style={box, fill=orange!10, draw=orange!60, text width=1.7cm},
    updatebox/.style={box, fill=red!8, draw=red!50, text width=2.6cm},
    membox/.style={box, fill=green!8, draw=green!50, text width=2.4cm},
    heurbox/.style={box, fill=gray!12, draw=gray!50, dashed, text width=1.8cm},
    arr/.style={-{Stealth[length=4pt]}, thick, draw=black!70},
    dasharr/.style={-{Stealth[length=4pt]}, thick, dashed, draw=black!40},
]

% Row 1: pipeline
\node[mainbox] (task) {Task\\Assignment};
\node[mainbox, right=of task] (appraisal) {Task Appraisal\\{\tiny (LLM hybrid)}};
\node[mainbox, right=of appraisal] (strategy) {Strategy\\Choice};
\node[mainbox, right=of strategy] (outcome) {Outcome\\{\tiny 6 types}};

% Row 2: classifiers
\node[classbox, below left=0.7cm and -0.3cm of outcome] (avoid) {Avoid\\分类\\{\tiny 3-class}};
\node[classbox, below right=0.7cm and -0.3cm of outcome] (support) {Support\\Mode\\{\tiny 2-class}};

% Row 3: update core
\node[updatebox, below=2.0cm of outcome] (update) {\texttt{apply\_helplessness\_update}\\{\tiny 成功路径 / 失败路径}};

% Mediators feeding into update
\node[mainbox, left=0.8cm of update, text width=1.7cm] (mediators) {SE\\Felt Control\\Uncontrollability};

% Protection and damping
\node[membox, right=0.8cm of update] (csm) {Controllable\\Success Memory\\{\tiny 乘性保护}};
\node[heurbox, below right=0.5cm and -0.5cm of update] (damp) {Damping\\{\tiny heuristic}};

% Row 4: memory update
\node[membox, below=1.5cm of update] (memupdate) {Experience Memory\\{\tiny SE / CSM / episode buffer}};

% Arrows: main pipeline
\draw[arr] (task) -- (appraisal);
\draw[arr] (appraisal) -- (strategy);
\draw[arr] (strategy) -- (outcome);

% Arrows: classifiers
\draw[arr] (outcome.south) -- ++(0,-0.3) -| (avoid.north);
\draw[arr] (outcome.south) -- ++(0,-0.3) -| (support.north);

% Arrows: into update
\draw[arr] (avoid.south) |- (update.west);
\draw[arr] (support.south) |- (update.east);
\draw[arr] (mediators.east) -- (update.west);
\draw[arr] (csm.west) -- (update.east);
\draw[dasharr] (damp.north west) -- (update.south east);

% Arrows: update -> memory
\draw[arr] (update.south) -- (memupdate.north);

% Feedback loop
\draw[dasharr] (memupdate.west) -- ++(-3.5,0) |- (strategy.south west);

% Appraisal provides mediators
\draw[arr] (appraisal.south) -- ++(0,-0.5) -| (mediators.north);

\end{tikzpicture}

\framebottomnote{虚线箭头表示跨事件反馈。灰色虚线框表示工程启发式（非文献直接推出）。}
\end{frame}

% ============================================================
\section{核心公式}
% ============================================================

\begin{frame}{核心公式}

\textbf{成功路径}（\texttt{success\_self} / \texttt{success\_with\_help}）

\vspace{0.15cm}
\[
H_{t+1} = \text{clamp}\!\bigl(H_t + \underbrace{\text{base}}_{\text{-1.4 / -0.7}} - \underbrace{\text{mastery\_recovery}}_{0 \sim 1.85}\,,\ 0,\ 100\bigr)
\]

\vspace{0.3cm}
\textbf{失败路径}（\texttt{failure} / \texttt{abandon} / \texttt{avoid}）

\vspace{0.15cm}
\[
\text{raw} = \underbrace{\text{base}}_{\substack{0.15 \\ \sim 0.9}}
+ \underbrace{\text{rep}}_{\substack{0 \\ \sim 1.5}}
+ \underbrace{\text{uncontroll}}_{\substack{0 \\ \sim 2.4}}
+ \underbrace{\text{eff\_loss}}_{\substack{0 \\ \sim 1.6}}
+ \underbrace{\text{ctrl\_loss}}_{\substack{0 \\ \sim 0.9}}
- \underbrace{\text{buffer}}_{\substack{0 \\ \sim 0.12}}
\]

\vspace{0.1cm}
\[
H_{t+1} = \text{clamp}\!\Bigl(H_t
+ \text{raw}
\times \underbrace{\alpha_{\text{avoid}}}_{\substack{0.15 \\ \sim 1.0}}
\times \underbrace{(1 - \text{CSM\_prot})}_{\substack{0.55 \\ \sim 1.0}}
\times \underbrace{\lambda(H_t)}_{\substack{0.55 \\ \sim 1.0}}
\,,\ 0,\ 100\Bigr)
\]

\end{frame}

% ============================================================
\section{机制--文献映射}
% ============================================================

\begin{frame}[t]{机制--文献映射（一）：理论主线}
\scriptsize
\setlength{\tabcolsep}{2.5pt}
\renewcommand{\arraystretch}{1.15}
\centering
\begin{tabular}{
    L{0.13\linewidth}
    L{0.24\linewidth}
    L{0.20\linewidth}
    L{0.33\linewidth}
}
\toprule
\textbf{机制} & \textbf{代码实现} & \textbf{理论依据} & \textbf{核心引句} \\
\midrule
不可控感\\主驱动
&
\texttt{UNCONTROLL\_DELTAS}$\{$1:1.2, 2:2.4$\}$远大于 \texttt{BASE\_DELTAS}$\{$fail:0.6$\}$
&
helplessness 核心是 action-outcome contingency 被破坏
&
Seligman (1967): ``Shock termination was independent of responding.'' \newline
Abramson (1978): ``Once people perceive noncontingency, they attribute their helplessness to a cause.'' \\

\midrule
SE 中介层
&
\texttt{efficacy\_loss} $= (50{-}\text{SE})/22$; SE 按 outcome $\times$ support\_mode 差异化更新
&
负面经历应先伤效能感，再间接影响 helplessness
&
Bandura (1977): ``Expectations of personal efficacy determine whether coping behavior will be initiated, how much effort will be expended, and how long it will be sustained.'' \\

\midrule
Felt control\\事件调节
&
\texttt{control\_loss} $= (50{-}\text{FC})/28$; FC 每轮由 LLM appraisal 生成
&
control 是多层构念，即时感受与跨事件 SE 应分开
&
Skinner (1996): ``The framework is used to analyze more than 100 terms, such as sense of control, proxy control, and primary control.'' \\
\bottomrule
\end{tabular}
\end{frame}

\begin{frame}[t]{机制--文献映射（二）：Avoid 拆分与可控成功记忆}
\scriptsize
\setlength{\tabcolsep}{2.5pt}
\renewcommand{\arraystretch}{1.15}
\centering
\begin{tabular}{
    L{0.13\linewidth}
    L{0.24\linewidth}
    L{0.20\linewidth}
    L{0.33\linewidth}
}
\toprule
\textbf{机制} & \textbf{代码实现} & \textbf{理论依据} & \textbf{核心引句} \\
\midrule
Avoid 拆分
&
helpless $\times$1.0 / risk $\times$0.35 / low-value $\times$0.15
&
不做 $\neq$ 无助；可能是风险规避或低价值判断
&
Davis (1989): ``usefulness had a significantly greater correlation with usage behavior than did ease of use.'' \newline
IRT review (2026): ``coded to the IRT domains: usage, value, risk, tradition, and image barriers.'' \\

\midrule
可控成功\\记忆
&
高质量 success 才积累 CSM（需 FC$\geq$45, U$\leq$1）；每天 $\times$0.985 衰减；乘性保护最高 45\%
&
先前可控经验应提供长期保护
&
Seligman (1967): ``Initial experience with escape ... prevented interference with later responding.'' \newline
Maier (2016): ``This passivity can be overcome by learning control.'' \newline
Review (2025): ``the sense of control offers the ability to effectively respond to environmental stressors.'' \\
\bottomrule
\end{tabular}
\end{frame}

\begin{frame}[t]{机制--文献映射（三）：Support 间接化与成功恢复}
\scriptsize
\setlength{\tabcolsep}{2.5pt}
\renewcommand{\arraystretch}{1.15}
\centering
\begin{tabular}{
    L{0.13\linewidth}
    L{0.24\linewidth}
    L{0.20\linewidth}
    L{0.33\linewidth}
}
\toprule
\textbf{机制} & \textbf{代码实现} & \textbf{理论依据} & \textbf{核心引句} \\
\midrule
Support\\间接修复
&
直接 buffer 极小（$\leq$0.12）；主效果经 SE 差异化更新（enabling +4.5 vs substituting +1.5）和 CSM 选择性积累传导
&
support 应先修复 control / efficacy，不是万能减分器
&
Digital feedback (2025): ``partially mediated by information processing self-efficacy.'' \newline
Digital experiences (2025): ``older persons may become more dependent on others.'' \\

\midrule
成功恢复\\差异化
&
self 最强恢复（$\Delta$max $\approx$ $-$3.25）；enabling help 次之（$\approx$ $-$1.40）；substituting 最弱（$\approx$ $-$1.20）
&
自己做成比被帮着做成更能恢复控制感
&
Bandura (1977): ``through experiences of mastery, further enhancement of self-efficacy and corresponding reductions in defensive behavior.'' \newline
Maier (2016): ``passivity can be overcome by learning control.'' \\
\bottomrule
\end{tabular}
\end{frame}

\begin{frame}[t]{机制--文献映射（四）：建模选择}
\scriptsize
\setlength{\tabcolsep}{2.5pt}
\renewcommand{\arraystretch}{1.15}
\centering
\begin{tabular}{
    L{0.13\linewidth}
    L{0.24\linewidth}
    L{0.20\linewidth}
    L{0.33\linewidth}
}
\toprule
\textbf{机制} & \textbf{代码实现} & \textbf{性质} & \textbf{与理论的关系} \\
\midrule
Damping
&
$\lambda(H) = 1 - 0.45 H/100$，clamp $[0.55, 1.0]$
&
\textbf{工程启发式}
&
与 Maier (2016)``被动是默认状态''方向相容，但不是文献直接推出的公式。主要目的是防止反馈环数值爆炸。 \\

\midrule
无自动回落
&
\texttt{helplessness\_score} 无 daily decay；只靠成功事件下拉
&
\textbf{建模选择}
&
与 Abramson (1978) ``learned attitude'' 方向相容：如果归因 stable，helplessness 不应自动消退。但具体``不加 decay''是操作化选择。 \\

\midrule
Trust /\\avoidance\\独立链
&
\texttt{trust\_in\_apps}、\texttt{avoidance\_tendency} 在 \texttt{compat.py} 中独立更新
&
\textbf{建模决策}
&
UTAUT (2003) 和 Skinner (1996) 支持``多维构念不应折叠''的一般原则，但具体三条独立链是工程选择。 \\
\bottomrule
\end{tabular}
\end{frame}

% ============================================================
\section{初步结果}
% ============================================================

\begin{frame}[t]{10-day 3-seed Results Snapshot}
\footnotesize
\begin{columns}[T]
\begin{column}{0.36\textwidth}
\textbf{核心观察}
\begin{itemize}
    \item \textbf{High Friction + Low Assist}：attempt 明显下降，negative share 与 helplessness 上升最大
    \item \textbf{High Friction + High Assist}：attempt 与 help-seeking 部分恢复，但 helplessness / trust 改善有限
    \item \textbf{Low Friction + High Assist}：五项主指标均向预期方向改善
\end{itemize}

\textbf{实验设置}
\begin{itemize}
    \item steady 单阶段
    \item 10 days，3 paired seeds
    \item 30-minute step，logical clock
\end{itemize}
\end{column}

\begin{column}{0.62\textwidth}
\scriptsize
\centering
\textbf{Paired effects vs baseline}\\[2pt]
\resizebox{\linewidth}{!}{
\begin{tabular}{lccccc}
\toprule
\textbf{Contrast vs baseline} & \textbf{$\Delta$Attempt} & \textbf{$\Delta$Neg.} & \textbf{$\Delta$Help.} & \textbf{$\Delta$Trust} & \textbf{$\Delta$Avoid.} \\
\midrule
High Friction + Low Assist
& -0.349
& +0.449
& +18.614
& -19.783
& +7.167 \\

High Friction + High Assist
& -0.116
& +0.195
& +14.645
& -22.017
& +12.033 \\

Low Friction + High Assist
& +0.109
& -0.170
& -27.158
& +40.083
& -35.408 \\
\bottomrule
\end{tabular}
}
\end{column}
\end{columns}

{\tiny Data source: \texttt{stage6\_single\_10day\_3seed\_logical\_clock\_20260408}. 所有 paired QC 通过；由于 $n{=}3$，此页主要用于展示 effect direction 与 magnitude。}
\end{frame}

% ============================================================
\section{总结}
% ============================================================

\begin{frame}{总结与下一步}

\begin{columns}[T]
\begin{column}{0.48\textwidth}
\textbf{已完成}
\begin{itemize}
    \item 10 条机制全部落地并逐条映射到文献
    \item 区分了三种证据层级：
    \begin{itemize}
        \scriptsize
        \item 文献直接支持（不可控感主驱动、SE 中介、avoid 拆分、CSM、support 间接化）
        \item 理论方向约束（无自动回落）
        \item 工程启发式（damping、独立更新链）
    \end{itemize}
    \item 核心公式（成功/失败两路径）完整呈现
\end{itemize}
\end{column}

\begin{column}{0.48\textwidth}
\textbf{下一步}
\begin{itemize}
    \item 补充 avoid 分类的 LLM hybrid 校准层
    \item 设计消融实验矩阵（6 组消融 + 5 项主指标）
    \item 跑固定 multi-seed baseline 并保存
    \item 完成 paired-seed 对照分析
    \item 整合为论文 Mechanism Section 的完整叙事
\end{itemize}
\end{column}
\end{columns}

\framebottomnote{消融矩阵：No-Uncontrollability / No-Efficacy-Mediator / No-Avoid-Split / No-Control-Memory / Direct-Support-Only / No-Damping。主指标：AttemptRate, NegShare, HelplessDelta, TrustDelta, AvoidanceDelta。}
\end{frame}

\begin{frame}
    \centering
    \Huge{请老师批评指正}
\end{frame}

\end{document}
